\section{Conclusion}
\label{sec:conclusion}

In disaggregated agentic serving, the cost of a discarded continuation is incurred at the prefill instance and the connector, whereas existing pruning signals are evaluated at the decode instance.
\sys is a cross-stage pruning framework: it applies pre-kernel admission from residency and a prefill-side heuristic score, defers uncertain candidates to a decode-side probe with a separately calibrated threshold, and uses a pruning-aware copy-on-write context tree to manage the KV state created by these decisions, retaining only the committed spine.
The accounting in \Cref{prop:acct} separates the savings that depend on the cache regime from those due to admission.
Empirically, prefill-side admission is complementary to three decode-time pruners and reduces computed prefill by $23\%$ at matched committed-answer accuracy; at $k{=}16$, a $128$-token probe increases coverage to $20/32$ at $0.69\times$ the full-fan-out wall-clock; and retaining only the committed spine reduces median peak KV usage to $0.72\times$ that of APC.

\section*{Limitations}
\label{sec:limitations}

Our scorers are lightweight heuristics evaluated on a controlled branching mix; natural agent residuals may be less separable, and learned scorers such as STOP~\cite{stoploss} can be substituted without changing the admission policy.
The admission experiments use moderate sample sizes ($n{=}16$ and $n{=}32$), so small differences in committed accuracy fall within run-to-run variation.
Extending the evaluation from reasoning and Tree-of-Thoughts sessions to tool-using agents with real environments is a natural next step.
